% Part: first-order-logic
% Chapter: syntax-and-semantics
% Section: free-vars-sentences

\documentclass[../../../include/open-logic-section]{subfiles}

\begin{document}

\olfileid{fol}{syn}{fvs}

\olsection{自由\printtoken{P}{variable}与\printtoken{P}{sentence}}

\begin{defn}[变元的自由出现]
\ollabel{defn:free-occ}
公式中变元的\emph{自由}出现归纳定义如下：
\begin{enumerate}
\item \indcase*{!A}{$!A$ 是原子公式}{$\indfrm$ 中所有变元的出现都是自由的。}

\tagitem{prvNot}{\indcase{!A}{\lnot !B}{$\indfrm$ 的自由变元出现恰好就是 $!B$ 的自由变元出现。}}{}

\item \indcase{!A}{(!B \ast !C)}{$\indfrm$ 的自由变元出现就是 $!B$ 与 $!C$ 中的自由变元出现。}

\tagitem{prvAll}{\indcase{!A}{\lforall[x][!B]}{$\indfrm$ 中的自由变元出现是 $!B$ 中除 $x$ 的出现之外的所有变元出现。}}{}

\tagitem{prvEx}{\indcase{!A}{\lexists[x][!B]}{$\indfrm$ 中的自由变元出现是 $!B$ 中除 $x$ 的出现之外的所有变元出现。}}{}
\end{enumerate}
\end{defn}

\begin{defn}[约束变元]
如果变元在公式~$!A$ 中的一个出现不是自由的，则该出现是\emph{约束}的。
\end{defn}

\begin{prob}
参照 \olref[fol][syn][fvs]{defn:free-occ} 的思路，给出约束变元出现的归纳定义。
\end{prob}

\begin{defn}[辖域]
\iftag{prvAll}{如果 $\lforall[x][!B]$ 是公式~$!A$ 中某子公式的出现，那么 $!B$ 在~$!A$ 中的相应出现就称为该 $\lforall[x]$ 出现的\emph{辖域}。\iftag{prvEx}{对 $\lexists[x]$ 同理。}{}}{如果 $\lexists[x][!B]$ 是公式~$!A$ 中某子公式的出现，那么 $!B$ 在~$!A$ 中的相应出现就称为该 $\lexists[x]$ 出现的\emph{辖域}。}

如果 $!B$ 是公式~$!A$ 中量词出现 \iftag{prvAll}{$\lforall[x]$\iftag{prvEx}{ 或 $\lexists[x]$}{}}{$\lexists[x]$} 的辖域，那么 $!B$ 中 $x$ 的自由出现在 \iftag{prvAll}{$\lforall[x][!B]$\iftag{prvEx}{ 和 $\lexists[x][!B]$}{}}{$\lexists[x][!B]$} 中则是约束的。我们称这些出现\emph{受}上述量词出现\emph{约束}。
\end{defn}

\begin{ex}
考虑以下公式：
\[
\lexists[\Obj v_0][\underbrace{\Atom{\Obj A^2_0}{\Obj v_0,\Obj v_1}}_{!B}] 
\]
$!B$ 代表 $\lexists[\Obj v_0]$ 的辖域。
该量词约束了 $!B$ 中 $\Obj v_0$ 的出现，但不约束 $\Obj v_1$ 的出现。因此，$\Obj v_1$ 在此情况下是自由变元。

我们现在可以看看这在更复杂的公式~$!A$ 中是如何起作用的：
\[
\lforall[\Obj v_0][\underbrace{(\Atom{\Obj A^1_0}{\Obj v_0} \lif
    \Atom{\Obj A^2_0}{\Obj v_0, \Obj v_1})}_{!B}] \lif \lexists[\Obj
  v_1][\underbrace{(\Atom{\Obj A^2_1}{\Obj v_0, \Obj v_1} \lor \lforall[\Obj v_0][\overbrace{\lnot \Atom{\Obj A^1_1}{\Obj v_0}}^{!D}])}_{!C}]
\]
$!B$ 是第一个 $\lforall[\Obj v_0]$ 的辖域，$!C$ 是 $\lexists[\Obj v_1]$ 的辖域，而 $!D$ 是第二个 $\lforall[\Obj v_0]$ 的辖域。第一个 $\lforall[\Obj v_0]$ 约束了~$!B$ 中 $\Obj v_0$ 的出现，$\lexists[\Obj v_1]$ 约束了 $!C$ 中 $\Obj v_1$ 的出现，第二个 $\lforall[\Obj v_0]$ 约束了~$!D$ 中 $\Obj v_0$ 的出现。$\Obj v_1$ 的第一次出现和 $\Obj v_0$ 的第四次出现在~$!A$ 中是自由的。$\Obj v_0$ 的最后一次出现在 $!D$ 中是自由的，但在 $!C$ 和~$!A$ 中则是约束的。
\end{ex}

\begin{defn}[语句]
一个公式~$!A$ 是\article{sentence} \emph{语句}，当且仅当它不含变元的自由出现。
\end{defn}

% add examples!

\end{document}